% ============================================================
% CAPITOLO 13 — Release: Pubblicazione sugli Store
% ============================================================
\chapter{Release: Pubblicazione sugli Store}\index{release}\index{Google Play}\index{App Store}

\section*{Cosa Costruirai}
In questo capitolo preparerai l'app Flutter per la pubblicazione:
\begin{itemize}
  \item Icona dell'app e splash screen generati dall'IA
  \item Build di produzione per Android (APK/AAB) e iOS (IPA)
  \item Firma digitale dell'app
  \item Listing sugli store (testi, screenshot)
  \item Processo di pubblicazione su Google Play e App Store
\end{itemize}

\textbf{Tempo stimato}: 90--120 minuti\\
\textbf{Prerequisito}: App Flutter funzionante (Cap.~11--12)

% ---------------------------------------------------------
\section{Preparazione alla Produzione}

\begin{praticabox}[title={Pratica --- Aggiorna il \texttt{\_CONTEXT.md} per il rilascio}]
Aggiungi una sezione al contesto:
\begin{lstlisting}[language={}]
## Release

### Identita App
- Nome: Notes App
- Bundle ID Android: com.example.notesapp
- Bundle ID iOS: com.example.notesapp
- Versione: 1.0.0+1

### Ambienti
- Dev: http://10.0.2.2:3000/api (emulatore)
- Staging: https://staging-api.notes-app.example.com/api
- Produzione: https://api.notes-app.example.com/api

### Build
- Android: flutter build appbundle (per Play Store)
- iOS: flutter build ipa (per App Store)
- Firma: keystore in android/app/keystore.jks (NON committare)
\end{lstlisting}
\end{praticabox}

% ---------------------------------------------------------
\section{Icona e Splash Screen}

\begin{praticabox}[title={Pratica --- Genera l'icona con l'IA}]
\begin{lstlisting}[language={}]
Genera un'icona per l'app Notes:
- Dimensione: 1024x1024 pixel
- Sfondo: gradiente indigo (da #4F46E5 a #3730A3)
- Icona centrale: nota stilizzata in bianco
- Stile: flat, minimal, bordi arrotondati
- Formato: PNG con trasparenza

Salva come assets/icon/app_icon.png
\end{lstlisting}
\end{praticabox}

\begin{tipbox}
Se l'IA non può generare immagini direttamente, chiedi di generare un'icona SVG con il codice, oppure usa strumenti come DALL-E o Midjourney con un prompt che l'IA ti aiuta a scrivere.
\end{tipbox}

\begin{praticabox}[title={Pratica --- Configura flutter\_launcher\_icons}]
\begin{lstlisting}[language={}]
1. Aggiungi la dipendenza dev in pubspec.yaml
2. Aggiungi la configurazione:
   flutter_launcher_icons:
     android: true
     ios: true
     image_path: "assets/icon/app_icon.png"
     adaptive_icon_background: "#4F46E5"
     adaptive_icon_foreground: "assets/icon/app_icon_foreground.png"
3. Genera le icone con: dart run flutter_launcher_icons
\end{lstlisting}
\end{praticabox}

\begin{praticabox}[title={Pratica --- Splash Screen}]
\begin{lstlisting}[language={}]
1. Aggiungi flutter_native_splash in pubspec.yaml
2. Configura:
   flutter_native_splash:
     color: "#4F46E5"
     image: assets/icon/app_icon.png
     android_12:
       color: "#4F46E5"
       icon_background_color: "#4F46E5"
       image: assets/icon/app_icon.png
3. Genera con: dart run flutter_native_splash:create
\end{lstlisting}
\end{praticabox}

% ---------------------------------------------------------
\section{Configurazione di Produzione}

\begin{praticabox}[title={Pratica --- Ambiente di produzione}]
\begin{lstlisting}[language={}]
Modifica api_config.dart per supportare ambienti:
1. Usa --dart-define per passare l'ambiente al build
2. In produzione, URL API punta al backend deployato
3. Disabilita log di debug in produzione
4. Configura: flutter build appbundle --dart-define=ENV=production
\end{lstlisting}
\end{praticabox}

L'IA dovrebbe generare qualcosa come:

\begin{lstlisting}[language=Java,title={ApiConfig con ambienti}]
class ApiConfig {
  static const String environment = String.fromEnvironment(
    'ENV',
    defaultValue: 'development',
  );

  static String get baseUrl {
    switch (environment) {
      case 'production':
        return 'https://api.notes-app.example.com/api';
      case 'staging':
        return 'https://staging-api.notes-app.example.com/api';
      default:
        return 'http://10.0.2.2:3000/api';
    }
  }
}
\end{lstlisting}

% ---------------------------------------------------------
\section{Build Android}

\subsection{Generare il Keystore}

Il keystore è il certificato digitale che firma la tua app. Senza di esso, Google Play non la accetta.

\begin{lstlisting}[language=bash]
keytool -genkey -v -keystore android/app/keystore.jks \
  -keyalg RSA -keysize 2048 -validity 10000 \
  -alias notes-app
\end{lstlisting}

\begin{warnbox}
Il file \texttt{keystore.jks} è \textsc{segreto}. \textsc{Non} committarlo in git. Se lo perdi, non potrai mai più aggiornare la tua app sullo store. Fai un backup sicuro.
\end{warnbox}

\begin{praticabox}[title={Pratica --- Configura la firma}]
\begin{lstlisting}[language={}]
1. Crea android/key.properties con i dati del keystore
   (storePassword, keyPassword, keyAlias, storeFile)
2. Modifica android/app/build.gradle per leggere key.properties
3. Aggiungi key.properties e keystore.jks al .gitignore
4. Configura: minifyEnabled true, shrinkResources true
\end{lstlisting}
\end{praticabox}

\subsection{Build}

\begin{lstlisting}[language=bash]
# App Bundle (per Google Play -- consigliato)
flutter build appbundle --dart-define=ENV=production

# APK (per distribuzione diretta)
flutter build apk --release --dart-define=ENV=production
\end{lstlisting}

Il file generato si trova in:
\begin{itemize}
  \item \textbf{AAB}: \texttt{build/app/outputs/bundle/release/app-release.aab}
  \item \textbf{APK}: \texttt{build/app/outputs/flutter-apk/app-release.apk}
\end{itemize}

% ---------------------------------------------------------
\section{Build iOS}

\begin{warnbox}
Per compilare per iOS hai bisogno di un Mac con Xcode installato e un account Apple Developer (\$99/anno).
\end{warnbox}

\begin{praticabox}[title={Pratica --- Configurazione Xcode}]
\begin{lstlisting}[language={}]
1. Apri ios/Runner.xcworkspace in Xcode
2. In Signing & Capabilities:
   - Seleziona il tuo Team (Apple Developer account)
   - Verifica il Bundle Identifier
3. In General:
   - Imposta la versione (1.0.0)
   - Imposta il build number (1)
4. In Build Settings:
   - iOS Deployment Target: 15.0 (o superiore)
\end{lstlisting}
\end{praticabox}

\begin{lstlisting}[language=bash]
flutter build ipa --dart-define=ENV=production
\end{lstlisting}

Il file IPA si trova in \texttt{build/ios/ipa/}.

% ---------------------------------------------------------
\section{Listing sugli Store}

\begin{praticabox}[title={Pratica --- Genera i testi con l'IA}]
\begin{lstlisting}[language={}]
Genera i testi per la pubblicazione sugli store:
1. Titolo: massimo 30 caratteri
2. Descrizione breve: massimo 80 caratteri
3. Descrizione completa: 4000 caratteri max
   - Paragrafo introduttivo
   - Lista funzionalita principali
   - Sezione sicurezza/privacy
   - Call to action finale
4. Parole chiave: 7-10 keywords rilevanti
5. Categoria: Produttivita
6. Content rating: Everyone / 4+

Genera sia la versione italiana che inglese.
\end{lstlisting}
\end{praticabox}

Per gli screenshot, esegui l'app sull'emulatore e cattura le schermate principali:

\begin{enumerate}
  \item \textbf{Login screen} --- con i bottoni OAuth
  \item \textbf{Dashboard} --- con alcune note di esempio
  \item \textbf{Dettaglio nota} --- la nota aperta
  \item \textbf{Creazione nota} --- il form di inserimento
  \item \textbf{Ricerca} --- risultati di una ricerca
\end{enumerate}

\begin{tipbox}
Usa \texttt{flutter screenshot} per catturare screenshot dall'emulatore direttamente dal terminale.
\end{tipbox}

% ---------------------------------------------------------
\section{Pubblicazione su Google Play}

\begin{enumerate}
  \item Vai su Google Play Console
  \item Crea una nuova app
  \item Compila la scheda del listing (testi, screenshot, icona)
  \item Vai in Release $\rightarrow$ Production $\rightarrow$ Create new release
  \item Carica il file \texttt{.aab}
  \item Compila la dichiarazione di conformità dei contenuti
  \item Imposta prezzo (gratuita) e paesi di distribuzione
  \item Invia per la revisione
\end{enumerate}

La prima revisione richiede generalmente 1--3 giorni lavorativi.

\subsection{Privacy Policy}

Google Play richiede una privacy policy. Chiedi all'IA:

\begin{praticabox}[title={Pratica --- Privacy policy}]
\begin{lstlisting}[language={}]
Genera una privacy policy per l'app Notes che:
- Spiega che l'app usa OAuth per l'autenticazione
- I dati delle note sono salvati su server sicuri
- Non condivide dati con terze parti
- L'utente puo eliminare account e tutti i dati
- Conforme al GDPR
- URL: https://notes-app.example.com/privacy
\end{lstlisting}
\end{praticabox}

% ---------------------------------------------------------
\section{Pubblicazione su App Store}

\begin{enumerate}
  \item Vai su App Store Connect
  \item Crea una nuova app
  \item Compila le informazioni (testi, screenshot, icona)
  \item In Xcode: Product $\rightarrow$ Archive $\rightarrow$ Distribute App $\rightarrow$ App Store Connect
  \item Torna su App Store Connect, seleziona la build
  \item Invia per la revisione
\end{enumerate}

La revisione Apple richiede generalmente 1--2 giorni lavorativi.

\begin{warnbox}
Apple è più restrittiva di Google. Assicurati che:
\begin{itemize}
  \item L'app funzioni senza crash
  \item Tutti i link (privacy policy, termini) siano raggiungibili
  \item L'app abbia un valore per l'utente (non sia solo una ``demo'')
  \item Il login OAuth funzioni correttamente
\end{itemize}
\end{warnbox}

% ---------------------------------------------------------
\section{CI/CD per Mobile}

\begin{praticabox}[title={Pratica --- GitHub Actions per build automatiche}]
\begin{lstlisting}[language={}]
Crea un workflow GitHub Actions che:
1. Si attiva su push al branch main
2. Fa il build Flutter per Android
3. Esegue flutter analyze e flutter test
4. Genera l'APK di release
5. Carica l'APK come artifact

File: .github/workflows/flutter-build.yml
\end{lstlisting}
\end{praticabox}

L'IA dovrebbe generare qualcosa come:

\begin{lstlisting}[language={},title={flutter-build.yml}]
name: Flutter Build

on:
  push:
    branches: [main]
  pull_request:
    branches: [main]

jobs:
  build:
    runs-on: ubuntu-latest
    steps:
      - uses: actions/checkout@v4
      - uses: subosito/flutter-action@v2
        with:
          flutter-version: '3.x'
          channel: 'stable'
      
      - name: Install dependencies
        working-directory: notes_mobile
        run: flutter pub get
      
      - name: Analyze
        working-directory: notes_mobile
        run: flutter analyze
      
      - name: Test
        working-directory: notes_mobile
        run: flutter test
      
      - name: Build APK
        working-directory: notes_mobile
        run: flutter build apk --release
             --dart-define=ENV=production
      
      - name: Upload APK
        uses: actions/upload-artifact@v4
        with:
          name: release-apk
          path: notes_mobile/build/app/outputs/
                flutter-apk/app-release.apk
\end{lstlisting}

\begin{notebox}
Per la pubblicazione automatica su Play Store, puoi aggiungere il plugin Fastlane al workflow. Ma per il primo rilascio, la pubblicazione manuale è sufficiente.
\end{notebox}

% ---------------------------------------------------------
\section*{Cosa Hai Costruito --- Progetto Finale}

In 13~capitoli, hai costruito un \textbf{sistema completo}:

\begin{figure}[H]
\centering
\includegraphics[width=0.85\textwidth]{06-notes-fullstack-arch}
\caption{Architettura finale Notes Full-Stack}
\end{figure}

% ---------------------------------------------------------
\section*{Riepilogo}

\begin{center}
\begin{tabularx}{\textwidth}{l X}
\toprule
\textbf{Aspetto} & \textbf{Dettaglio} \\
\midrule
Icona & flutter\_launcher\_icons con icona adattiva \\
Splash & flutter\_native\_splash con tema indigo \\
Build Android & AAB firmato con keystore \\
Build iOS & IPA via Xcode (richiede Mac + Apple Developer) \\
CI/CD & GitHub Actions: analyze + test + build \\
Store listing & Testi generati dall'IA in IT e EN \\
Privacy & Policy GDPR generata dall'IA \\
\bottomrule
\end{tabularx}
\end{center}
